% =================================================================
% sprint1_trung.tex
% Sprint 1 — Theory Lead: Đào Việt Trung
% Nội dung: Mục 1.1 (đầy đủ) và Mục 1.2 (trừ 1.2.3)
%
% Cách dùng:
% (1) Compile standalone: pdflatex sprint1_trung.tex (hoặc xelatex)
% (2) Tích hợp vào repo của Hiếu: copy phần body (từ \chapter ... đến trước \end{document})
%     vào file chương 1 của report chính.
% =================================================================

\documentclass[12pt,a4paper]{report}

% --- Vietnamese support ---
% File này được viết để compile bằng XeLaTeX hoặc LuaLaTeX
% (lệnh: xelatex sprint1_trung.tex)
% Nếu Hiếu dùng pdflatex + vntex trong repo chính, thay 3 dòng dưới bằng:
%   \usepackage[utf8]{inputenc}
%   \usepackage[T5]{fontenc}
%   \usepackage[vietnamese]{babel}
\usepackage{fontspec}
\usepackage{polyglossia}
\setmainlanguage{vietnamese}
\setotherlanguage{english}

% --- Toán học ---
\usepackage{amsmath,amsthm,amssymb,amsfonts}
\usepackage{mathtools}

% --- Thuật toán ---
\usepackage{algorithm}
\usepackage{algpseudocode}
\floatname{algorithm}{Thuật toán}
\renewcommand{\algorithmicrequire}{\textbf{Đầu vào:}}
\renewcommand{\algorithmicensure}{\textbf{Đầu ra:}}
\renewcommand{\algorithmicwhile}{\textbf{khi}}
\renewcommand{\algorithmicdo}{\textbf{thì}}
\renewcommand{\algorithmicif}{\textbf{nếu}}
\renewcommand{\algorithmicthen}{\textbf{thì}}
\renewcommand{\algorithmicreturn}{\textbf{trả về}}

% --- Định dạng ---
\usepackage{enumitem}
\usepackage{booktabs}
\usepackage{array}
\usepackage{geometry}
\geometry{margin=2.5cm}
\usepackage{hyperref}
\hypersetup{colorlinks=true,linkcolor=blue!50!black,citecolor=blue!50!black}
\usepackage{xcolor}

% --- Theorem environments ---
\newtheorem{theorem}{Định lý}[chapter]
\newtheorem{lemma}[theorem]{Bổ đề}
\newtheorem{corollary}[theorem]{Hệ quả}
\newtheorem{proposition}[theorem]{Mệnh đề}
\theoremstyle{definition}
\newtheorem{definition}[theorem]{Định nghĩa}
\newtheorem{example}[theorem]{Ví dụ}
\theoremstyle{remark}
\newtheorem{remark}[theorem]{Nhận xét}

% --- Title ---
\title{Sprint 1 — Nền tảng lý thuyết\\
\large Mục 1.1 và 1.2 (trừ 1.2.3)}
\author{Đào Việt Trung \\ \small Theory Lead}
\date{}

\begin{document}
\maketitle

\chapter{Kiến thức cơ sở}
\label{chap:foundations}

Chương này trình bày các kiến thức nền tảng từ lý thuyết số, đại số trừu tượng và độ phức tạp tính toán cần thiết để hiểu lược đồ chữ ký số ElGamal. Các khái niệm được xây dựng theo trình tự tăng dần: từ phép toán cơ bản trên số nguyên (Mục~\ref{sec:number-theory}), qua cấu trúc nhóm và các tính chất khai thác được trong mật mã, đến bài toán logarit rời rạc — bài toán \emph{khó} mà toàn bộ độ an toàn của ElGamal dựa vào (Mục~\ref{sec:dlp}).

% ===================================================================
\section{Lý thuyết số và đại số trừu tượng}
\label{sec:number-theory}

\subsection{Số nguyên tố và phép đồng dư}
\label{subsec:primes-modular}

\subsubsection{Số nguyên tố và định lý cơ bản của số học}

\begin{definition}[Số nguyên tố]
Số nguyên tố là số nguyên dương $p > 1$ chỉ có đúng hai ước dương là $1$ và chính nó. Tập các số nguyên tố ký hiệu là $\mathbb{P}$.
\end{definition}

\begin{theorem}[Euclid]
Tập các số nguyên tố là vô hạn.
\end{theorem}

\begin{theorem}[Định lý cơ bản của số học]
Mọi số nguyên $n > 1$ đều có duy nhất một cách viết dưới dạng
\[
n = p_1^{e_1} p_2^{e_2} \cdots p_r^{e_r},
\]
trong đó $p_1 < p_2 < \cdots < p_r$ là các số nguyên tố và $e_i \in \mathbb{N}^*$.
\end{theorem}

Trong mật mã, ta đặc biệt quan tâm đến các số nguyên tố \emph{lớn} — thường có độ dài từ $2048$ bit trở lên — vì chúng tạo nên nền tảng cho các bài toán khó như phân tích thừa số nguyên (RSA) và logarit rời rạc (ElGamal, DSA, ECDSA).

\subsubsection{Phép chia có dư và phép đồng dư}

\begin{theorem}[Phép chia có dư]
Với mọi $a \in \mathbb{Z}$ và $n \in \mathbb{Z}^+$, tồn tại duy nhất cặp số nguyên $(q, r)$ sao cho
\[
a = qn + r, \quad 0 \leq r < n.
\]
Ta gọi $q$ là \emph{thương}, $r$ là \emph{số dư}, ký hiệu $r = a \bmod n$.
\end{theorem}

\begin{definition}[Đồng dư]
Hai số nguyên $a, b$ được gọi là \emph{đồng dư} modulo $n$, ký hiệu
\[
a \equiv b \pmod{n},
\]
nếu $n \mid (a - b)$, tức $a$ và $b$ có cùng số dư khi chia cho $n$.
\end{definition}

\begin{proposition}[Tính chất cơ bản của đồng dư]
Phép đồng dư là một quan hệ tương đương trên $\mathbb{Z}$ và tương thích với các phép toán cộng, trừ, nhân. Cụ thể, nếu $a \equiv b \pmod{n}$ và $c \equiv d \pmod{n}$ thì:
\begin{enumerate}
    \item $a \pm c \equiv b \pm d \pmod{n}$,
    \item $ac \equiv bd \pmod{n}$,
    \item $a^k \equiv b^k \pmod{n}$ với mọi $k \in \mathbb{N}$.
\end{enumerate}
\end{proposition}

Tính chất này cho phép thực hiện toàn bộ các phép tính của ElGamal trong tập hữu hạn $\{0, 1, \ldots, p-1\}$ — không bao giờ bị tràn số dù $p$ có cỡ $2048$ bit hay lớn hơn.

Tập các lớp tương đương modulo $n$ cùng phép cộng và phép nhân tạo thành \emph{vành thương} $\mathbb{Z}/n\mathbb{Z}$, ký hiệu tắt là $\mathbb{Z}_n$. Tập các phần tử khả nghịch (đối với phép nhân) của $\mathbb{Z}_n$ ký hiệu là $\mathbb{Z}_n^*$.

\subsubsection{Định lý Fermat nhỏ và định lý Euler}

\begin{theorem}[Fermat nhỏ]
\label{thm:fermat}
Cho $p$ là số nguyên tố và $a \in \mathbb{Z}$ với $\gcd(a, p) = 1$. Khi đó
\[
a^{p-1} \equiv 1 \pmod{p}.
\]
Hệ quả: với mọi $a \in \mathbb{Z}$, ta có $a^{p} \equiv a \pmod{p}$.
\end{theorem}

\begin{definition}[Hàm Euler $\varphi$]
Hàm Euler $\varphi: \mathbb{Z}^+ \to \mathbb{Z}^+$ đếm số các số nguyên trong $\{1, 2, \ldots, n\}$ nguyên tố cùng nhau với $n$:
\[
\varphi(n) = |\{a : 1 \leq a \leq n,\ \gcd(a, n) = 1\}|.
\]
Đặc biệt, $\varphi(p) = p - 1$ khi $p$ là số nguyên tố.
\end{definition}

\begin{theorem}[Euler — tổng quát hóa Fermat nhỏ]
Với mọi $n \in \mathbb{Z}^+$ và $a \in \mathbb{Z}$ thỏa $\gcd(a, n) = 1$:
\[
a^{\varphi(n)} \equiv 1 \pmod{n}.
\]
\end{theorem}

Định lý Fermat nhỏ đóng vai trò then chốt trong bộ công cụ mật mã:
\begin{itemize}[leftmargin=*]
    \item Cho phép tính nghịch đảo nhanh \emph{khi modulo là số nguyên tố}: $a^{-1} \equiv a^{p-2} \pmod{p}$.
    \item Là cơ sở cho các thuật toán kiểm tra tính nguyên tố xác suất như Fermat test, Miller--Rabin --- được dùng để sinh các số nguyên tố lớn trong giai đoạn KeyGen của ElGamal.
    \item Là nền tảng để chứng minh tính đúng đắn của lược đồ chữ ký ElGamal: lũy thừa trên $\mathbb{Z}_p^*$ chỉ phụ thuộc vào số mũ modulo $p-1$.
\end{itemize}

\subsubsection{Nghịch đảo modular và thuật toán Euclid mở rộng}

\begin{definition}[Nghịch đảo modular]
Nghịch đảo modular của $a$ theo modulo $n$ là số nguyên $a^{-1}$ thỏa mãn
\[
a \cdot a^{-1} \equiv 1 \pmod{n}.
\]
\end{definition}

\begin{theorem}
Nghịch đảo modular của $a$ modulo $n$ tồn tại khi và chỉ khi $\gcd(a, n) = 1$.
\end{theorem}

Thuật toán Euclid mở rộng (Extended Euclidean Algorithm) cho phép tính đồng thời $\gcd(a, n)$ và các \emph{hệ số Bézout} $(u, v) \in \mathbb{Z}^2$ thỏa mãn
\[
au + nv = \gcd(a, n).
\]
Khi $\gcd(a, n) = 1$, ta có $au \equiv 1 \pmod{n}$, do đó $a^{-1} \equiv u \pmod{n}$.

\begin{algorithm}[H]
\caption{Thuật toán Euclid mở rộng}
\label{alg:extgcd}
\begin{algorithmic}[1]
\Require $a, n \in \mathbb{Z}^+$
\Ensure $(d, u, v)$ với $d = \gcd(a, n)$ và $au + nv = d$
\State $(r_0, r_1) \gets (a, n)$
\State $(u_0, u_1) \gets (1, 0)$
\State $(v_0, v_1) \gets (0, 1)$
\While{$r_1 \neq 0$}
    \State $q \gets \lfloor r_0 / r_1 \rfloor$
    \State $(r_0, r_1) \gets (r_1, r_0 - q \cdot r_1)$
    \State $(u_0, u_1) \gets (u_1, u_0 - q \cdot u_1)$
    \State $(v_0, v_1) \gets (v_1, v_0 - q \cdot v_1)$
\EndWhile
\State \Return $(r_0, u_0, v_0)$
\end{algorithmic}
\end{algorithm}

Độ phức tạp: $O(\log \min(a, n))$ phép chia, mỗi phép chia trên số $k$-bit tốn $O(k^2)$ bit-operation. Tổng cộng $O(k^3)$ với $k = \log n$. Trong Python: hàm \texttt{pow(a, -1, n)} (Python 3.8+) đã tích hợp sẵn.

\begin{example}
Tính $7^{-1} \bmod 26$. Áp dụng thuật toán Euclid trên $26$ và $7$:
\begin{align*}
26 &= 3 \cdot 7 + 5, \\
7  &= 1 \cdot 5 + 2, \\
5  &= 2 \cdot 2 + 1, \\
2  &= 2 \cdot 1 + 0.
\end{align*}
Quay ngược các hệ số Bézout:
\begin{align*}
1 &= 5 - 2\cdot 2 \\
  &= 5 - 2(7 - 5) = 3\cdot 5 - 2\cdot 7 \\
  &= 3(26 - 3\cdot 7) - 2\cdot 7 = 3\cdot 26 - 11\cdot 7.
\end{align*}
Vậy $7 \cdot (-11) \equiv 1 \pmod{26}$, hay $7^{-1} \equiv 15 \pmod{26}$.
\end{example}

\begin{remark}
\label{rmk:euclid-elgamal}
Trong lược đồ ElGamal, ta cần tính $k^{-1} \bmod (p-1)$ ở bước ký. Lưu ý modulo ở đây là $p-1$ --- \emph{không nguyên tố} (vì $p-1$ luôn chẵn nên có thừa số $2$). Do đó \textbf{bắt buộc} phải dùng thuật toán Euclid mở rộng, không thể dùng công thức $a^{p-2}$ của định lý Fermat nhỏ. Đây là chi tiết cài đặt dễ bị nhầm khi mới làm quen với ElGamal.
\end{remark}

\subsubsection{Định lý số dư Trung Hoa}

\begin{theorem}[Định lý số dư Trung Hoa --- CRT]
Cho $n_1, n_2, \ldots, n_k$ là các số nguyên dương đôi một nguyên tố cùng nhau. Đặt $N = n_1 n_2 \cdots n_k$. Khi đó với mọi bộ số $(a_1, a_2, \ldots, a_k)$, hệ phương trình đồng dư
\[
x \equiv a_i \pmod{n_i}, \quad i = 1, 2, \ldots, k
\]
có nghiệm duy nhất modulo $N$.
\end{theorem}

CRT cho phép phân rã các phép tính trên $\mathbb{Z}_N$ thành các phép tính song song trên các $\mathbb{Z}_{n_i}$ nhỏ hơn. Đây cũng là công cụ then chốt trong thuật toán Pohlig--Hellman tấn công DLP khi $p-1$ có cấu trúc thuận lợi --- điều giải thích vì sao phải chọn $p$ là \emph{safe prime} (xem Mục~\ref{subsubsec:safe-prime}).

% --------------------------------------------------------------
\subsection{Nhóm hữu hạn và phần tử nguyên thủy}
\label{subsec:groups}

\subsubsection{Khái niệm nhóm và nhóm con}

\begin{definition}[Nhóm]
Một \emph{nhóm} là cặp $(G, \cdot)$ gồm tập $G \neq \emptyset$ và phép toán hai ngôi $\cdot : G \times G \to G$ thỏa mãn:
\begin{enumerate}
    \item \textbf{Kết hợp}: $(a \cdot b) \cdot c = a \cdot (b \cdot c)$ với mọi $a, b, c \in G$.
    \item \textbf{Có phần tử đơn vị}: tồn tại $e \in G$ sao cho $e \cdot a = a \cdot e = a$ với mọi $a \in G$.
    \item \textbf{Có phần tử nghịch đảo}: với mỗi $a \in G$, tồn tại $a^{-1} \in G$ sao cho $a \cdot a^{-1} = a^{-1} \cdot a = e$.
\end{enumerate}
Nhóm $(G, \cdot)$ được gọi là \emph{Abel} (hay \emph{giao hoán}) nếu thêm điều kiện $a \cdot b = b \cdot a$ với mọi $a, b \in G$. \emph{Cấp} của nhóm, ký hiệu $|G|$, là số phần tử của $G$.
\end{definition}

\begin{definition}[Nhóm con]
Tập con $H \subseteq G$ được gọi là \emph{nhóm con} của $G$, ký hiệu $H \leq G$, nếu $H$ cùng phép toán hạn chế từ $G$ là một nhóm. Tương đương: $H$ đóng đối với phép toán và phép lấy nghịch đảo.
\end{definition}

\begin{theorem}[Lagrange]
\label{thm:lagrange}
Nếu $H$ là nhóm con của nhóm hữu hạn $G$ thì $|H|$ là ước của $|G|$.
\end{theorem}

Trong mật mã khóa công khai dựa trên DLP, ta luôn làm việc với các nhóm \emph{hữu hạn, Abel}.

\subsubsection{Cấp của phần tử và nhóm cyclic}

\begin{definition}[Cấp của phần tử]
Cho nhóm $(G, \cdot)$ và $a \in G$. \emph{Cấp} của $a$, ký hiệu $\mathrm{ord}(a)$, là số nguyên dương nhỏ nhất $k$ sao cho $a^k = e$. Nếu không tồn tại $k$ như vậy, ta nói $a$ có cấp vô hạn.
\end{definition}

\begin{corollary}
Trong nhóm hữu hạn, $\mathrm{ord}(a) \mid |G|$ với mọi $a \in G$.
\end{corollary}

\begin{proof}
Tập $\langle a \rangle = \{a^0, a^1, \ldots, a^{\mathrm{ord}(a)-1}\}$ là một nhóm con của $G$, với cấp đúng bằng $\mathrm{ord}(a)$. Áp dụng định lý Lagrange.
\end{proof}

\begin{definition}[Nhóm cyclic]
Nhóm $G$ được gọi là \emph{cyclic} nếu tồn tại phần tử $g \in G$ sao cho mọi phần tử của $G$ đều có dạng $g^k$ với $k \in \mathbb{Z}$. Phần tử $g$ được gọi là \emph{phần tử sinh} (generator) của $G$, và ta viết $G = \langle g \rangle$.
\end{definition}

\begin{theorem}[Cấu trúc nhóm cyclic]
Mọi nhóm cyclic hữu hạn cấp $n$ đều đẳng cấu với $(\mathbb{Z}_n, +)$. Với mỗi ước số $d$ của $n$, nhóm cyclic cấp $n$ có đúng \emph{một} nhóm con cấp $d$, và nhóm con này cũng cyclic.
\end{theorem}

\subsubsection{Nhóm nhân \texorpdfstring{$\mathbb{Z}_p^*$}{Zp*}}

Cho $p$ là số nguyên tố. Tập
\[
\mathbb{Z}_p^* = \{1, 2, \ldots, p-1\}
\]
cùng với phép nhân modulo $p$ tạo thành một nhóm. Đây chính là cấu trúc đại số mà ElGamal hoạt động trên đó.

\begin{theorem}[Gauss]
\label{thm:zp-cyclic}
Với $p$ nguyên tố, $\mathbb{Z}_p^*$ là nhóm cyclic cấp $p - 1$.
\end{theorem}

Hệ quả của định lý này là tồn tại ít nhất một phần tử $g \in \mathbb{Z}_p^*$ sao cho mọi phần tử khác đều là một lũy thừa của $g$:
\[
\mathbb{Z}_p^* = \{g^0, g^1, g^2, \ldots, g^{p-2}\}.
\]

\subsubsection{Phần tử nguyên thủy}

\begin{definition}[Phần tử nguyên thủy]
Phần tử $g \in \mathbb{Z}_p^*$ được gọi là \emph{phần tử nguyên thủy} (hay \emph{căn nguyên thủy}, primitive root, generator) modulo $p$ nếu $\mathrm{ord}(g) = p - 1$, tức $g$ sinh ra toàn bộ $\mathbb{Z}_p^*$.
\end{definition}

\begin{theorem}[Số lượng phần tử nguyên thủy]
Số phần tử nguyên thủy modulo $p$ bằng $\varphi(p-1)$.
\end{theorem}

\paragraph{Cách kiểm tra phần tử nguyên thủy.}
Cho $p - 1 = q_1^{e_1} q_2^{e_2} \cdots q_r^{e_r}$ là phân tích thừa số nguyên tố. Phần tử $g \in \mathbb{Z}_p^*$ là nguyên thủy modulo $p$ khi và chỉ khi với mọi $i = 1, \ldots, r$:
\[
g^{(p-1)/q_i} \not\equiv 1 \pmod{p}.
\]

Lý do: nếu $g^{(p-1)/q_i} \equiv 1$ với một $i$ nào đó, thì $\mathrm{ord}(g)$ chia hết $(p-1)/q_i$, tức $\mathrm{ord}(g) < p-1$. Ngược lại, nếu mọi điều kiện trên đều thỏa, thì $\mathrm{ord}(g)$ không thể là ước thực sự của $p-1$ (vì mọi ước thực sự đều chia hết một $(p-1)/q_i$ nào đó), do đó $\mathrm{ord}(g) = p-1$.

\begin{example}[Phần tử nguyên thủy modulo $11$]
\label{ex:primroot-11}
Với $p = 11$, ta có $p - 1 = 10 = 2 \cdot 5$. Kiểm tra $g = 2$:
\begin{itemize}
    \item $2^{10/2} = 2^5 = 32 \equiv 10 \pmod{11} \neq 1$ \quad ✓
    \item $2^{10/5} = 2^2 = 4 \pmod{11} \neq 1$ \quad ✓
\end{itemize}
Vậy $g = 2$ là phần tử nguyên thủy của $\mathbb{Z}_{11}^*$. Liệt kê các lũy thừa:
\[
\begin{array}{c|cccccccccc}
i & 1 & 2 & 3 & 4 & 5 & 6 & 7 & 8 & 9 & 10 \\
\hline
2^i \bmod 11 & 2 & 4 & 8 & 5 & 10 & 9 & 7 & 3 & 6 & 1
\end{array}
\]
Sinh đủ $10$ phần tử của $\mathbb{Z}_{11}^*$. Số phần tử nguyên thủy modulo $11$ là $\varphi(10) = 4$, cụ thể là $\{2, 6, 7, 8\}$.
\end{example}

\subsubsection{Safe prime}
\label{subsubsec:safe-prime}

\begin{definition}[Safe prime]
Số nguyên tố $p$ được gọi là \emph{safe prime} nếu $q = (p-1)/2$ cũng là số nguyên tố. Khi đó $q$ được gọi là \emph{số nguyên tố Sophie Germain}.
\end{definition}

Trong thực hành mật mã, ta luôn chọn $p$ là safe prime vì ba lý do quan trọng:

\begin{enumerate}[leftmargin=*]
    \item \textbf{Kiểm tra phần tử nguyên thủy cực nhanh.} Phân tích $p - 1 = 2 \cdot q$ chỉ có hai thừa số, nên chỉ cần kiểm tra hai điều kiện $g^{(p-1)/2} \neq 1$ và $g^{(p-1)/q} = g^2 \neq 1$.
    
    \item \textbf{Hầu như mọi phần tử đều có cấp lớn.} Theo định lý Lagrange, cấp của phần tử là ước của $p-1 = 2q$, do đó chỉ có thể là $1, 2, q$, hoặc $2q$. Chỉ có $g = 1$ có cấp $1$ và $g = p-1$ có cấp $2$. Mọi phần tử khác đều có cấp $q$ hoặc $2q$ --- tức cùng cỡ với $p$.
    
    \item \textbf{Chống tấn công Pohlig--Hellman.} Tấn công này tận dụng các thừa số nhỏ của $p - 1$ để giải DLP nhanh trên từng nhóm con cấp nhỏ rồi ghép lại bằng CRT. Khi $p - 1 = 2q$ với $q$ nguyên tố lớn, không tồn tại nhóm con nhỏ để khai thác.
\end{enumerate}

% --------------------------------------------------------------
\subsection{Lũy thừa modular}
\label{subsec:modexp}

Bài toán tính $g^x \bmod p$ với $g, x, p$ lớn (cỡ vài nghìn bit) là phép toán cơ bản nhất trong mọi hệ mã khóa công khai dựa trên DLP. Nó xuất hiện ở cả ba bước KeyGen, Sign và Verify của ElGamal.

\subsubsection{Tại sao cần thuật toán hiệu quả}

Tính trực tiếp $g^x$ rồi mới rút gọn modulo $p$ là \emph{không khả thi}: với $g$ và $x$ cỡ $2048$ bit, giá trị $g^x$ có thể chứa tới $2^{2048} \cdot 2048 \approx 10^{619}$ chữ số --- vượt xa khả năng lưu trữ của mọi máy tính.

Ý tưởng khắc phục là rút gọn modulo $p$ \emph{sau mỗi phép nhân}, giữ cho các giá trị trung gian luôn nằm trong khoảng $[0, p^2)$. Đồng thời, ta dùng biểu diễn nhị phân của số mũ để giảm số phép nhân từ $x$ xuống chỉ còn $O(\log x)$.

\subsubsection{Thuật toán bình phương và nhân}

\paragraph{Ý tưởng.} Biểu diễn $x$ dưới dạng nhị phân $x = (b_{n-1} b_{n-2} \ldots b_1 b_0)_2$. Khi đó
\[
g^x = g^{\sum_{i=0}^{n-1} b_i 2^i} = \prod_{i:\, b_i = 1} g^{2^i}.
\]
Các giá trị $g^{2^i}$ được tính lần lượt bằng bình phương: $g, g^2, g^4, g^8, \ldots$, mỗi bước rút gọn modulo $p$.

\begin{algorithm}[H]
\caption{Lũy thừa modular --- bình phương và nhân}
\label{alg:modexp}
\begin{algorithmic}[1]
\Require $g, x, p$ với $x \geq 0$, $p > 0$
\Ensure $g^x \bmod p$
\State $\mathit{result} \gets 1$
\State $\mathit{base} \gets g \bmod p$
\While{$x > 0$}
    \If{$x$ lẻ}
        \State $\mathit{result} \gets \mathit{result} \cdot \mathit{base} \bmod p$
    \EndIf
    \State $x \gets \lfloor x / 2 \rfloor$
    \State $\mathit{base} \gets \mathit{base}^2 \bmod p$
\EndWhile
\State \Return $\mathit{result}$
\end{algorithmic}
\end{algorithm}

\paragraph{Phân tích độ phức tạp.}
Mỗi vòng lặp \texttt{while} luôn thực hiện một phép bình phương. Phép nhân ở dòng~5 chỉ thực hiện khi bit hiện tại của $x$ là $1$. Số vòng lặp đúng bằng $\lfloor \log_2 x \rfloor + 1$, do đó tổng số phép nhân không vượt quá $2 \log_2 x$.

Với $x$ cỡ $2048$ bit, ta chỉ cần khoảng $4000$ phép nhân số lớn --- hoàn toàn thực hiện được trong vài mili giây trên phần cứng phổ thông. Tổng độ phức tạp bit-operation là $O(\log^3 p)$ với phép nhân thông thường, có thể giảm xuống $O(\log^2 p \log\log p)$ với thuật toán Karatsuba hoặc FFT.

Trong Python, hàm tích hợp \texttt{pow(g, x, p)} đã cài đặt thuật toán này, kèm các tối ưu như Montgomery multiplication và sliding window.

\begin{example}
\label{ex:modexp}
Tính $3^{13} \bmod 17$. Ta có $13 = (1101)_2 = 8 + 4 + 1$, do đó
\[
3^{13} = 3^{8} \cdot 3^{4} \cdot 3^{1}.
\]
Bình phương liên tiếp modulo $17$:
\begin{align*}
3^1 &= 3, \\
3^2 &= 9, \\
3^4 &= 81 \equiv 13 \pmod{17}, \\
3^8 &\equiv 13^2 = 169 \equiv 16 \pmod{17}.
\end{align*}
Suy ra
\[
3^{13} \equiv 16 \cdot 13 \cdot 3 = 624 \equiv 12 \pmod{17}.
\]
Kiểm chứng nhanh: $3^{16} \equiv 1 \pmod{17}$ (Fermat nhỏ), và $3 \cdot 12 = 36 \equiv 2 \pmod{17}$, $3^{14} \equiv 2 \pmod{17}$, $3^{15} \equiv 6$, $3^{16} \equiv 18 \equiv 1$ ✓.
\end{example}

\subsubsection{Tính bất đối xứng --- nền tảng cho mật mã khóa công khai}

Phép lũy thừa modular có một tính chất kỳ diệu mà mật mã khóa công khai khai thác triệt để:

\begin{center}
\renewcommand{\arraystretch}{1.3}
\begin{tabular}{ll}
\toprule
\textbf{Chiều thuận} & $x \longmapsto y = g^x \bmod p$ là \emph{dễ}: $O(\log^3 p)$ bit-op. \\
\textbf{Chiều ngược} & $(g, y, p) \longmapsto x$ (DLP) là \emph{khó}: dưới hàm mũ. \\
\bottomrule
\end{tabular}
\end{center}

Tính bất đối xứng "dễ tính theo một chiều, khó tính theo chiều ngược" này chính là tính chất \emph{cửa sập một chiều} (one-way trapdoor) --- nền tảng triết lý của mật mã khóa công khai do Diffie và Hellman đề xuất năm 1976. Mục~\ref{sec:dlp} tiếp theo sẽ đi sâu vào bài toán logarit rời rạc và phân tích lý do cộng đồng mật mã tin tưởng nó là khó.

% ===================================================================
\section{Bài toán Logarit rời rạc}
\label{sec:dlp}

Mục này phát biểu bài toán logarit rời rạc (Discrete Logarithm Problem --- DLP), thảo luận tại sao DLP được tin tưởng là khó về mặt tính toán, và liên hệ với một số lớp bài toán trong lý thuyết độ phức tạp.

\subsection{Phát biểu bài toán}
\label{subsec:dlp-statement}

\subsubsection{Logarit thông thường và động cơ}

Trong giải tích thực, hàm $\log_g(\cdot)$ ngược với phép lũy thừa $g^{(\cdot)}$: nếu $y = g^x$ thì $x = \log_g y$. Phép lấy logarit thực là \emph{dễ} (có công thức và máy tính giải tích sẵn).

Khi chuyển sang nhóm hữu hạn cyclic, định nghĩa logarit vẫn được giữ nguyên về mặt cấu trúc, nhưng tính chất \emph{dễ tính} thì biến mất hoàn toàn --- đây chính là điểm mấu chốt mà mật mã khai thác.

\subsubsection{Phát biểu trong nhóm cyclic}

\begin{definition}[Bài toán logarit rời rạc --- DLP]
\label{def:dlp}
Cho $(G, \cdot)$ là nhóm cyclic hữu hạn cấp $n$, sinh bởi phần tử $g$. Cho $y \in G$. Tìm số nguyên $x \in \{0, 1, \ldots, n-1\}$ sao cho
\[
g^x = y.
\]
Số $x$ này được gọi là \emph{logarit rời rạc cơ số $g$} của $y$ trong $G$, ký hiệu $x = \log_g y$.
\end{definition}

\begin{proposition}[Tính tồn tại và duy nhất]
Vì $G = \langle g \rangle$ cyclic cấp $n$, mọi phần tử $y \in G$ đều có biểu diễn duy nhất dạng $g^x$ với $x \in \{0, 1, \ldots, n-1\}$. Do đó logarit rời rạc luôn tồn tại và duy nhất trong khoảng $[0, n-1]$.
\end{proposition}

\subsubsection{Phát biểu cụ thể trên \texorpdfstring{$\mathbb{Z}_p^*$}{Zp*}}

Cụ thể hóa Định nghĩa~\ref{def:dlp} cho nhóm $\mathbb{Z}_p^*$ với $p$ nguyên tố và $g$ là phần tử nguyên thủy:

\begin{quote}
\textbf{DLP trên $\mathbb{Z}_p^*$.} Cho số nguyên tố $p$, phần tử nguyên thủy $g \in \mathbb{Z}_p^*$, và $y \in \mathbb{Z}_p^*$. Tìm số nguyên $x \in \{0, 1, \ldots, p-2\}$ sao cho
\[
g^x \equiv y \pmod{p}.
\]
\end{quote}

Đây chính là dạng DLP mà lược đồ chữ ký ElGamal dựa vào.

\begin{example}[DLP với tham số nhỏ]
Với $p = 11, g = 2$ (đã chứng minh ở Ví dụ~\ref{ex:primroot-11} là phần tử nguyên thủy). Cho $y = 7$, tìm $x$ sao cho $2^x \equiv 7 \pmod{11}$.

Tra bảng các lũy thừa của $2$: $2^7 \equiv 7 \pmod{11}$, do đó $x = 7$. Vậy $\log_2 7 = 7$ trong $\mathbb{Z}_{11}^*$.

\smallskip
Khi $p$ nhỏ, bài toán này tầm thường (vét cạn). Khi $p \approx 2^{2048}$, không gian tìm kiếm lên tới $2^{2048}$ khả năng --- vượt xa mọi máy tính cổ điển.
\end{example}

\subsubsection{Tính chất một chiều}

DLP có cấu trúc đặc trưng cho \emph{hàm một chiều} (one-way function):
\begin{itemize}[leftmargin=*]
    \item Tính theo chiều thuận $x \mapsto g^x \bmod p$ là \emph{dễ} --- đã thấy ở Mục~\ref{subsec:modexp} với độ phức tạp $O(\log^3 p)$.
    \item Tính theo chiều ngược $y \mapsto \log_g y$ là \emph{khó} --- chưa có thuật toán đa thức nào được biết.
\end{itemize}

Quan trọng: tính một chiều của lũy thừa modular \emph{không phải} là điều hiển nhiên về mặt toán học. Trong giải tích thực, logarit \emph{dễ tính} không kém gì lũy thừa --- chính cấu trúc \emph{rời rạc} của $\mathbb{Z}_p^*$ tạo ra độ khó. Một cách trực giác: phép lũy thừa modular "trộn" thứ tự các phần tử trong $\mathbb{Z}_p^*$ theo cách rất "ngẫu nhiên", không tuân theo bất kỳ quy luật giải tích nào có thể đảo ngược.

\subsection{Tại sao DLP khó}
\label{subsec:dlp-hardness}

\subsubsection{Không có thuật toán đa thức được biết}

Tính đến nay (2026), \textbf{không có thuật toán đa thức nào được biết để giải DLP trên $\mathbb{Z}_p^*$ với $p$ tổng quát}. Các thuật toán tốt nhất hiện hành có độ phức tạp:
\begin{itemize}[leftmargin=*]
    \item \emph{Tổng quát} (áp dụng cho mọi nhóm cyclic): $O(\sqrt{p})$ thời gian --- thuật toán Pollard rho, Baby-step Giant-step.
    \item \emph{Riêng cho $\mathbb{Z}_p^*$} (khai thác cấu trúc số học): \emph{dưới hàm mũ} (sub-exponential), cụ thể là $L_p[1/3, c]$ --- thuật toán Index Calculus và các biến thể như Number Field Sieve.
\end{itemize}
Tất cả các thuật toán trên đều có độ phức tạp lớn hơn đa thức theo $\log p$. Với $p$ cỡ $2048$ bit, kể cả thuật toán tốt nhất cũng cần thời gian tính tương đương phân tích thừa số một số RSA-$2048$ bit --- vượt khả năng tính toán hiện tại.

Niềm tin về độ khó của DLP đến từ:
\begin{enumerate}[leftmargin=*]
    \item \emph{Lịch sử nghiên cứu lâu dài}: bài toán đã được nghiên cứu hơn $50$ năm trong cộng đồng mật mã và lý thuyết số, chưa có đột phá lớn nào ngoài Index Calculus (đã có từ những năm 1970).
    \item \emph{Đầu tư thực tế}: các tổ chức như NSA, NIST liên tục đánh giá lại độ an toàn của DLP và đến nay vẫn khuyến nghị nó là cơ sở cho các chuẩn quan trọng (DSA, DH).
    \item \emph{Tính phổ quát}: cùng bài toán xuất hiện tự nhiên trong nhiều lĩnh vực toán học, độ khó không phụ thuộc vào một ứng dụng cụ thể nào.
\end{enumerate}

\subsubsection{Liên hệ với lý thuyết độ phức tạp}

Để đặt DLP vào ngữ cảnh lý thuyết phức tạp, ta nhắc lại nhanh các lớp bài toán quyết định cơ bản:
\begin{itemize}[leftmargin=*]
    \item $\mathbf{P}$: lớp các bài toán quyết định giải được trong thời gian đa thức.
    \item $\mathbf{NP}$: lớp các bài toán có chứng nhận đáp án "có" \emph{kiểm chứng được} trong thời gian đa thức.
    \item $\mathbf{coNP}$: lớp các bài toán có chứng nhận đáp án "không" kiểm chứng được trong thời gian đa thức.
\end{itemize}

Phiên bản quyết định của DLP có dạng:
\begin{quote}
\textsc{DLP-Decision}: Cho $(p, g, y, k)$, có tồn tại $x \in \{0, \ldots, k\}$ sao cho $g^x \equiv y \pmod{p}$ không?
\end{quote}

\begin{proposition}
\textsc{DLP-Decision} thuộc $\mathbf{NP} \cap \mathbf{coNP}$.
\end{proposition}

\begin{proof}[Phác thảo]
\emph{Thuộc $\mathbf{NP}$:} Chứng nhận cho câu trả lời "có" là chính giá trị $x$. Việc kiểm tra $g^x \equiv y \pmod{p}$ và $x \leq k$ làm được trong thời gian đa thức bằng thuật toán bình phương và nhân.

\emph{Thuộc $\mathbf{coNP}$:} Chứng nhận cho câu trả lời "không" là logarit rời rạc \emph{thực} $x_0 = \log_g y$, kèm chứng minh $x_0 > k$. Việc kiểm tra cũng làm được trong thời gian đa thức.
\end{proof}

\begin{remark}[Vị trí của DLP trong lý thuyết phức tạp]
DLP nằm trong $\mathbf{NP} \cap \mathbf{coNP}$ nhưng \emph{chưa được chứng minh} thuộc $\mathbf{P}$, cũng \emph{chưa được chứng minh} là $\mathbf{NP}$-khó. Đây là tình thế "vùng xám" mà DLP chia sẻ với bài toán phân tích thừa số nguyên (factoring). Hai đặc điểm:
\begin{itemize}[leftmargin=*]
    \item \emph{Không thể $\mathbf{NP}$-khó} (trừ khi $\mathbf{NP} = \mathbf{coNP}$, điều được tin là không xảy ra): mọi bài toán $\mathbf{NP}$-khó nằm trong $\mathbf{NP} \cap \mathbf{coNP}$ sẽ kéo theo $\mathbf{NP} = \mathbf{coNP}$.
    \item \emph{Có thể thuộc $\mathbf{P}$}: chưa ai loại trừ được khả năng tồn tại thuật toán đa thức cho DLP. Đây là rủi ro nền của mọi hệ mã dựa trên DLP --- nếu có ai tìm ra thuật toán đa thức, toàn bộ hệ thống sụp đổ.
\end{itemize}
Quan trọng: \emph{độ khó của DLP là một giả thiết, không phải định lý}. Mật mã khóa công khai về bản chất luôn xây dựng trên các giả thiết về độ khó tính toán chưa được chứng minh.
\end{remark}

\subsubsection{Liên hệ với các bài toán mật mã liên quan}

Quanh DLP có một họ các bài toán liên quan mật thiết, sắp xếp theo độ khó tăng dần:

\begin{description}[leftmargin=*]
    \item[\textsc{DLP} (Discrete Logarithm Problem).] Cho $(g, y)$, tìm $x = \log_g y$. Đây là bài toán "khó nhất" trong họ.
    
    \item[\textsc{CDH} (Computational Diffie--Hellman).] Cho $(g, g^a, g^b)$, tính $g^{ab}$. Nếu có thể giải DLP thì có thể giải CDH (tìm $a$ từ $g^a$, rồi tính $(g^b)^a$). Chiều ngược chưa biết.
    
    \item[\textsc{DDH} (Decisional Diffie--Hellman).] Cho $(g, g^a, g^b, z)$, quyết định xem $z = g^{ab}$ hay $z$ ngẫu nhiên. Đây là bài toán "dễ nhất" trong họ.
\end{description}

Ta luôn có chuỗi rút gọn (reduction):
\[
\text{Giải DLP} \implies \text{Giải CDH} \implies \text{Giải DDH}.
\]
Do đó nếu DLP khó thì CDH và DDH cũng khó. Lược đồ chữ ký ElGamal dựa trên DLP, lược đồ trao đổi khóa Diffie--Hellman dựa trên CDH, các hệ mã đối xứng kiểu ElGamal dựa trên DDH.

\subsubsection{Ảnh hưởng của kích thước tham số}

Độ khó của DLP tăng theo kích thước của $p$, nhưng \emph{không tăng tuyến tính}:

\begin{center}
\renewcommand{\arraystretch}{1.2}
\begin{tabular}{rll}
\toprule
\textbf{Bit của $p$} & \textbf{Đánh giá an toàn} & \textbf{Ghi chú} \\
\midrule
$512$ & Đã phá & Đã có ghi nhận phá thực tế, chỉ dùng demo \\
$1024$ & Yếu, không khuyến nghị & NIST deprecated từ 2013 \\
$2048$ & An toàn ngắn--trung hạn (đến $\sim 2030$) & Mức tối thiểu hiện hành \\
$3072$ & An toàn dài hạn & Khuyến nghị mới của NIST \\
$7680$ & Tương đương AES-256 & Hiếm dùng do tốc độ chậm \\
\bottomrule
\end{tabular}
\end{center}

Sự gia tăng kích thước $p$ kéo theo:
\begin{itemize}[leftmargin=*]
    \item \emph{Độ an toàn} tăng dưới hàm mũ theo Index Calculus.
    \item \emph{Thời gian ký/xác minh} tăng theo $O(\log^3 p)$.
\end{itemize}
Sự đánh đổi này là bài toán thực tiễn mà mục~$2.4$ (thực nghiệm) sẽ định lượng cụ thể.

\subsubsection{Đe dọa từ máy tính lượng tử}

Năm $1994$, Peter Shor công bố thuật toán lượng tử giải DLP và factoring trong thời gian đa thức. Cụ thể, trên máy tính lượng tử lý tưởng, DLP có thể giải trong $O(\log^3 p)$ phép tính lượng tử.

Hiện tại (2026), các máy tính lượng tử thực tế chưa đủ lớn để chạy Shor trên $p$ kích thước mật mã. Tuy nhiên đây là rủi ro \emph{dài hạn} đáng kể, là động lực cho nghiên cứu \emph{mật mã hậu lượng tử} (post-quantum cryptography) --- các hệ mã dựa trên các bài toán khó mới như lattice problems, code-based, hash-based, multivariate polynomial.

\bigskip

% ===================================================================
\noindent\hrulefill

\noindent\small\textit{Kết thúc Mục $1.1$ và $1.2$ (trừ $1.2.3$). Mục $1.2.3$ --- "Các thuật toán tấn công DLP đã biết" --- và các mục tiếp theo của Chương $1$ sẽ được trình bày trong các phần riêng.}

\end{document}
